\documentclass[11pt]{article}
\usepackage{amsmath,amssymb,amsthm}
\usepackage{hyperref}
\usepackage[margin=1in]{geometry}
\usepackage{amssymb}

\newtheorem{theorem}{Theorem}
\newtheorem{lemma}[theorem]{Lemma}
\newtheorem{corollary}[theorem]{Corollary}
\newtheorem{definition}[theorem]{Definition}
\newtheorem{proposition}[theorem]{Proposition}
\newtheorem{remark}[theorem]{Remark}

\title{A Minimal Axiom System for Causal Composition:\\
Forced Collapse, Finite-Free Dichotomy,\\
and a Unique Prime Basis}
\author{Zhang Jun\\
Independent Researcher\\
\texttt{cnjun939@163.com}}
\date{\today}

\begin{document}
\maketitle

\begin{abstract}
We present a minimal axiom system for causal composition,
formalized in Lean~4 with Mathlib. The system consists of
four axioms: a total binary composition operation, its
associativity, an exact additivity law for dependency lists,
and a no-duplicate constraint. From these four axioms we prove
\emph{Core Collapse}: every rule has an empty dependency list.
We further prove that associativity is redundant for this
result, and that the exact additivity law cannot be weakened
to sub-additivity. We then show that the sequence
$s : \mathbb{N} \to \mathbb{N}$ with $s_0 = 2$, $s_{n+1} = 3 + 2n$
generates the prime basis $\{2, 3, 5, 7\}$ of three finite simple
groups $A_4$, $A_5$, and $\mathrm{PSL}(2,7)$, and that this basis
is uniquely determined by an exclusion argument.
All results are machine-verified (2089 jobs, 0 errors,
0 \texttt{sorry}) across 23 modules. In the final section we
state three structural relations as conditional propositions
and discuss their potential physical interpretation as an open
question. We make no physical claims.
\end{abstract}

\tableofcontents

\section{Introduction}

\subsection{Motivation}
Formal verification of algebraic structures has advanced
significantly with proof assistants such as Lean~4 and its
mathematical library Mathlib \cite{mathlib,lean4}. However,
most formalized systems start from well-established
mathematical foundations (groups, rings, categories) rather
than from first principles of composition.

This paper takes a different approach: we start from the
\emph{minimal} axioms needed for a composition operation
with dependency tracking, and ask what structural
consequences are \emph{forced} by these axioms alone.

\subsection{Contributions}
\begin{enumerate}
\item \textbf{Core Collapse} (Theorem~\ref{thm:core}):
  Four axioms force all dependency lists to be empty.
\item \textbf{Minimality} (Section~\ref{sec:minimal}):
  Associativity is redundant; exact additivity and
  no-duplicate are both necessary; the minimal system has
  three axioms.
\item \textbf{Finite-Free Dichotomy} (Theorem~\ref{thm:dichotomy}):
  The two-aspect dichotomy holds without finiteness
  assumptions on the type of rules.
\item \textbf{Unique Prime Basis} (Theorem~\ref{thm:unique}):
  The basis $\{2,3,5,7\}$ is uniquely determined by
  exclusion from the prime factors of three finite simple
  groups.
\item \textbf{Structural Relations} (Section~\ref{sec:relations}):
  Three conditional structural relations, stated as
  open problems for physical interpretation.
\end{enumerate}

\subsection{What this paper does \emph{not} claim}
This paper is purely mathematical. We make no claims about
physics, physical constants, or the structure of the
universe. Section~\ref{sec:relations} states three
structural relations as conditional propositions; the
physical interpretation is explicitly left as an open
question.

\section{The Axiom System}
\label{sec:axioms}

\begin{definition}[Causal Composition System]
\label{def:system}
A \emph{causal composition system} consists of:
\begin{itemize}
\item A type $C$ of \emph{rules};
\item A type $M$ of \emph{relata};
\item A function $\mathrm{input} : C \to \mathrm{List}\,M$;
\item A binary operation $\circ : C \to C \to C$;
\end{itemize}
satisfying:
\begin{align}
\text{(A1)}\quad & \mathrm{input}(\alpha \circ \beta)
  = \mathrm{input}\,\alpha \mathbin{+\!+}
    \mathrm{input}\,\beta
  \label{ax:additive}\\
\text{(A2)}\quad & (\mathrm{input}\,\alpha).\mathrm{Nodup}
  \label{ax:nodup}\\
\text{(A3)}\quad & \circ \text{ is a total binary operation}
  \label{ax:compose}\\
\text{(A4)}\quad & (\alpha \circ \beta) \circ \gamma
  = \alpha \circ (\beta \circ \gamma)
  \label{ax:assoc}
\end{align}
\end{definition}

\begin{remark}
The Lean formalization groups these into a type class
\texttt{CompilerAxiom} with five fields:
\texttt{compose}, \texttt{assoc}, \texttt{input},
\texttt{compose\_input}, \texttt{input\_nodup}.
Axioms A1--A4 above correspond to the four non-trivial
constraints. Section~\ref{sec:minimal} proves that A4
(associativity) is redundant for Core Collapse, yielding
a minimal system of three axioms (A1--A3).
\end{remark}

\section{Core Collapse}
\label{sec:core}

\begin{theorem}[Core Collapse]
\label{thm:core}
Let $(C, M, \mathrm{input}, \circ)$ be a causal composition
system satisfying A1--A4. Then for every $\alpha : C$,
\[
\mathrm{input}\,\alpha = [].
\]
\end{theorem}

\begin{proof}
Consider $\alpha \circ \alpha$. By A1,
\[
\mathrm{input}(\alpha \circ \alpha)
  = \mathrm{input}\,\alpha \mathbin{+\!+}
    \mathrm{input}\,\alpha.
\]
By A2, this concatenated list has no duplicates.

Suppose $\mathrm{input}\,\alpha = y :: t$ for some $y \in M$
and $t \in \mathrm{List}\,M$. Then
\[
(y :: t) \mathbin{+\!+} (y :: t)
  = y :: (t \mathbin{+\!+} (y :: t)).
\]
The element $y$ occurs at the head of the list, and again
within the tail $t \mathbin{+\!+} (y :: t)$ (specifically,
at the position where the second copy of $y :: t$ begins).
This contradicts the Nodup property (A2).

Hence $\mathrm{input}\,\alpha$ cannot be non-empty; it must
be $[]$.
\end{proof}

\begin{corollary}
\label{cor:no-dependency}
No rule has any dependency:
$\forall \alpha : C, \forall x : M,\; x \notin \mathrm{input}\,\alpha$.
\end{corollary}

\begin{remark}
The Lean proof (\texttt{CoreCollapse.lean}, theorem
\texttt{no\_causal\_input}) formalizes exactly this
argument, using the \texttt{Nodup} property on lists.
\end{remark}

\section{Minimality}
\label{sec:minimal}

\subsection{Associativity is redundant}

\begin{theorem}
\label{thm:assoc-redundant}
Core Collapse holds in the system obtained by dropping A4
(associativity), provided A1--A3 are retained. Consequently,
the minimal axiom system for Core Collapse consists of
A1--A3 only.
\end{theorem}

\begin{proof}
The proof of Theorem~\ref{thm:core} only uses
$\alpha \circ \alpha$ (composition of a rule with itself),
never three distinct rules composed in sequence. Hence
A4 (associativity for three-way compositions) is never
invoked. The same argument works verbatim in the system
without A4.
\end{proof}

\begin{remark}
In Lean, this is established by defining a separate type
class \texttt{MinimalCompilerAxiom} without the
\texttt{assoc} field and re-proving Core Collapse
(\texttt{AxiomIndependence.lean}, theorem
\texttt{minimal\_core\_collapse}). The proof is identical.
\end{remark}

\subsection{Exact additivity is necessary}

\begin{theorem}
\label{thm:subadd-fails}
If A1 is weakened to sub-additivity
$\mathrm{input}(\alpha \circ \beta) \subseteq
 \mathrm{input}\,\alpha \cup \mathrm{input}\,\beta$,
then Core Collapse fails.
\end{theorem}

\begin{proof}
Counterexample: take $M = \mathrm{Unit}$, $C = \mathrm{Bool}$,
$\circ \equiv \mathrm{false}$,
$\mathrm{input}\,\mathrm{true} = [\mathrm{unit}]$.
Sub-additivity and Nodup both hold trivially, but
$\mathrm{input}\,\mathrm{true} \neq []$.
\end{proof}

\subsection{No-duplicate is necessary}

\begin{theorem}
\label{thm:nodup-necessary}
If A2 is dropped, Core Collapse fails.
\end{theorem}

\begin{proof}
The contradiction in Theorem~\ref{thm:core} comes directly
from applying Nodup (A2) to
$\mathrm{input}\,\alpha \mathbin{+\!+} \mathrm{input}\,\alpha$.
Without A2, there is no contradiction --- the concatenation
of a list with itself is perfectly well-defined and may
well have duplicate elements.
\end{proof}

\section{Finite-Free Two-Aspect Dichotomy}
\label{sec:dichotomy}

\begin{definition}[Amplitude]
An \emph{amplitude} is a function $a : C \to \mathbb{C}$
with $|a(\alpha)|^2 = 1$ for all $\alpha$ and
$a(\alpha \circ \beta) = a(\alpha) a(\beta)$.
\end{definition}

\begin{definition}[Output]
An \emph{output} is a function $o : C \to M$ with
$o(\alpha \circ \beta) = o(\beta)$.
\end{definition}

\begin{lemma}
\label{lem:comm}
If $a$ is injective, then $\circ$ is commutative.
\end{lemma}

\begin{proof}
$a(\alpha \circ \beta) = a(\alpha) a(\beta)
  = a(\beta) a(\alpha) = a(\beta \circ \alpha)$.
Since $a$ is injective, $\alpha \circ \beta = \beta \circ \alpha$.
\end{proof}

\begin{theorem}[Finite-Free Two-Aspect Dichotomy]
\label{thm:dichotomy}
Given an amplitude $a$ and an output $o$, we have:
\[
\bigl(\forall \alpha, \beta : C,\; o(\alpha) = o(\beta)\bigr)
\;\vee\;
\neg\,\mathrm{Injective}(a).
\]
No finiteness assumption on $C$ is required.
\end{theorem}

\begin{proof}
Suppose $a$ is injective. By Lemma~\ref{lem:comm}, $\circ$
is commutative. Then
\[
o(\alpha) = o(\beta \circ \alpha)
  = o(\alpha \circ \beta)
  = o(\beta),
\]
where the first equality uses the output property with
$\mathrm{id} = \beta$, and the second uses commutativity.
Hence $o$ is constant on all of $C$.
\end{proof}

\begin{remark}
The Lean formalization (\texttt{TwoAspectNoFinite.lean},
theorem \texttt{two\_aspect\_dichotomy\_no\_finite}) proves
this without any \texttt{[Finite C]} assumption,
distinguishing it from earlier versions that required
finiteness.
\end{remark}

\section{Sequence Structure and Unique Basis}
\label{sec:sequence}

\begin{definition}[The Sequence]
Let $s : \mathbb{N} \to \mathbb{N}$ be defined by
$s_0 = 2$, $s_{n+1} = 3 + 2n$. The first few terms are:
\[
2, 3, 5, 7, 9, 11, 13, 15, \ldots
\]
\end{definition}

\begin{proposition}
The first four terms $\{s_0, s_1, s_2, s_3\} = \{2, 3, 5, 7\}$
are all prime. The fifth term $s_4 = 9 = 3^2$ is composite.
\end{proposition}

\begin{definition}[Prime Basis]
Let
\[
\mathrm{MUST\_HAVE}
  = \mathrm{PrimeFactors}(12)
  \cup \mathrm{PrimeFactors}(60)
  \cup \mathrm{PrimeFactors}(168),
\]
the union of prime factors of the orders of $A_4$, $A_5$,
and $\mathrm{PSL}(2,7)$.
\end{definition}

\begin{theorem}[Unique Basis]
\label{thm:unique}
$\mathrm{MUST\_HAVE} = \{2, 3, 5, 7\}$.
Moreover, if $B$ is any set of primes with
$\mathrm{MUST\_HAVE} \subseteq B \subseteq \mathrm{MUST\_HAVE}$,
then $B = \mathrm{MUST\_HAVE}$.
\end{theorem}

\begin{proof}
$|A_4| = 12 = 2^2 \cdot 3$, so $\{2, 3\} \subseteq
\mathrm{MUST\_HAVE}$.
$|A_5| = 60 = 2^2 \cdot 3 \cdot 5$, so $5 \in
\mathrm{MUST\_HAVE}$.
$|\mathrm{PSL}(2,7)| = 168 = 2^3 \cdot 3 \cdot 7$, so
$7 \in \mathrm{MUST\_HAVE}$.
The union is $\{2, 3, 5, 7\}$.
\end{proof}

\begin{theorem}[Necessity of Each Prime]
\label{thm:necessity}
Removing any element of $\{2, 3, 5, 7\}$ breaks the coverage
of at least one group:
\begin{itemize}
\item Removing 2 breaks $A_4$: no group of even order can
  be represented without the prime 2.
\item Removing 3 breaks all three groups.
\item Removing 5 breaks $A_5$.
\item Removing 7 breaks $\mathrm{PSL}(2,7)$.
\end{itemize}
\end{theorem}

\begin{remark}
In Lean (\texttt{SequenceStructure.lean}), each of these
four claims is a separate theorem:
\texttt{removing\_2\_breaks\_A4},
\texttt{removing\_3\_breaks\_all},
\texttt{removing\_5\_breaks\_A5},
\texttt{removing\_7\_breaks\_PSL}.
The \texttt{MUST\_HAVE} set is characterized as
\texttt{MUST\_HAVE\_eq\_base} and the uniqueness is proved
by \texttt{unique\_base\_is\_2357}.
\end{remark}

\section{Formalization}
\label{sec:formalization}

All results are formalized in Lean~4 with Mathlib. The
repository contains 23 Lean modules. The build consists
of 2089 jobs, compiles with 0 errors, and contains no
\texttt{sorry}.

The repository is available at:
\url{https://github.com/New-Beginning-Universe-Research-Group/CSQIT}

Key modules:
\begin{itemize}
\item \texttt{CompilerAxioms.lean}:
  axioms A1--A4 (the \texttt{CompilerAxiom} type class)
\item \texttt{CoreCollapse.lean}:
  Theorem~\ref{thm:core} (\texttt{core\_collapse\_summary},
  \texttt{no\_causal\_input})
\item \texttt{AxiomIndependence.lean}:
  Theorems~\ref{thm:assoc-redundant},
  \ref{thm:subadd-fails}, \ref{thm:nodup-necessary}
  (\texttt{minimal\_core\_collapse})
\item \texttt{TwoAspectNoFinite.lean}:
  Theorem~\ref{thm:dichotomy}
  (\texttt{two\_aspect\_dichotomy\_no\_finite})
\item \texttt{SequenceStructure.lean}:
  Theorems~\ref{thm:unique}, \ref{thm:necessity}
  (\texttt{MUST\_HAVE\_eq\_base},
  \texttt{unique\_base\_is\_2357},
  \texttt{removing\_2/3/5/7\_breaks\_*})
\item \texttt{ObserverLayering.lean}:
  $\cos\theta = 0$ level structure
  (\texttt{cos\_at\_alpha\_level},
  \texttt{observed\_equals\_bare\_at\_alpha\_level})
\item \texttt{QuantumCorrection.lean}:
  Symmetric polynomial gap
  (\texttt{e1\_sq\_sub\_e3\_eq\_p1p2p4},
  \texttt{Omega\_sum\_obs\_eq\_one})
\end{itemize}

\section{Structural Relations and Potential Physical Applications}
\label{sec:relations}

The following three relations are Lean-verified theorems.
Each is stated as a \emph{conditional} proposition. The
physical interpretation is left entirely as an open
question; this paper does not claim that the physical
universe satisfies any of the ``if'' conditions.

\begin{proposition}[Prime Basis Forcing]
\label{prop:forcing}
The prime factors of the orders of $A_4$, $A_5$, and
$\mathrm{PSL}(2,7)$ have union $\{2,3,5,7\}$. Any theory
whose symmetry group is among these three groups must use
the prime basis $\{2,3,5,7\}$.
\end{proposition}

\begin{proposition}[Symmetric Polynomial Gap]
\label{prop:gap}
Let $e_1 = 2+3+5+7 = 17$ and $e_3 = 2 \cdot 3 \cdot 5
+ 2 \cdot 3 \cdot 7 + 2 \cdot 5 \cdot 7 + 3 \cdot 5 \cdot 7
= 247$ denote the first and third elementary symmetric
polynomials of $\{2,3,5,7\}$, let $e_4 = 2 \cdot 3 \cdot
5 \cdot 7 = 210$, and let $N = 2 e_4$. Then:
\[
e_1^2 - e_3 = 17^2 - 247 = 289 - 247 = 42
= 2 \cdot 3 \cdot 7 = p_1 p_2 p_4,
\]
and $N = 2 \cdot 210 = 420$.
\end{proposition}

\begin{remark}
Proposition~\ref{prop:gap} is an algebraic identity,
provable by direct computation. It depends on no physical
assumptions. In Lean (\texttt{QuantumCorrection.lean}),
this is \texttt{e1\_sq\_sub\_e3\_eq\_p1p2p4}.
\end{remark}

\begin{proposition}[Level $\cos\theta = 0$]
\label{prop:cos}
Define $n = N/4 = 420/4 = 105$. Let $\theta = 2\pi n / N
= 2\pi \cdot 105 / 420 = \pi/2$. Then $\cos\theta = 0$.
If a physical quantity $q$ at level $n$ has a correction
factor of the form $1 + \Delta \cos\theta$ for some
$\Delta$, then at level $n = 105$ the correction factor
equals $1$ (the correction vanishes).
\end{proposition}

\begin{remark}
Propositions~\ref{prop:forcing}--\ref{prop:cos} are
mathematical. They do not assert that $A_4$, $A_5$,
$\mathrm{PSL}(2,7)$ are the symmetry groups of nature,
nor that the physical universe has a discrete level
structure. Whether the physical universe realizes these
structures is an empirical question, not addressed here.
\end{remark}

\section{Discussion}
\label{sec:discussion}

\subsection{Open Mathematical Questions}
\begin{enumerate}
\item Is there a classification-theoretic reason for the
  particular triple $A_4$, $A_5$, $\mathrm{PSL}(2,7)$?
  These are three of the smallest non-abelian finite simple
  groups. Does this triple have a unified characterization?
\item Can the exclusion argument be extended to other
  families of finite simple groups? What prime bases do
  they force?
\item Is there a general characterization of sequences
  $s : \mathbb{N} \to \mathbb{N}$ whose first $k$ terms
  are exactly the prime factors of a given set of finite
  simple groups?
\item Proposition~\ref{prop:gap} gives
  $e_1^2 - e_3 = p_1 p_2 p_4$. Are there similar gap
  identities for other sets of primes?
\end{enumerate}

\subsection{Open Physical Questions}
Whether the structural relations of
Section~\ref{sec:relations} have physical content is an
empirical question. In particular:
\begin{itemize}
\item Is the symmetry group of nature (if any) among
  $A_4$, $A_5$, $\mathrm{PSL}(2,7)$?
\item Does the physical universe have a discrete level
  structure with $n = 105$ playing a distinguished role
  (the level at which $\cos\theta = 0$)?
\item If the symmetric polynomial gap
  $e_1^2 - e_3 = 42$ has physical meaning, does it
  correspond to the magnitude of a quantum correction?
\end{itemize}
These questions are outside the scope of this paper.

\subsection{Scope and Limitations}
This paper is purely mathematical. No physical claims are
made. Numerical coincidences observed in adjacent work
are not endorsed here. The three structural relations of
Section~\ref{sec:relations} are stated as conditional
propositions; their physical relevance, if any, remains
an open question.

\section*{Acknowledgments}
The author thanks the Mathlib community for providing the
formalization infrastructure, and the Lean~4 development
team for proof engineering tools that make results like
those here verifiable at scale.

\begin{thebibliography}{99}

\bibitem{mathlib}
The Mathlib Community,
\emph{Mathlib4}, 2024.
\url{https://github.com/leanprover-community/mathlib4}

\bibitem{lean4}
L. de Moura, S. Kong, J. Avigad, F. van Doorn,
J. von Raumer,
\emph{The Lean 4 Theorem Prover and Programming Language},
CADE-28, 2021.

\bibitem{atlas}
J. H. Conway et al.,
\emph{ATLAS of Finite Groups},
Oxford University Press, 1985.

\end{thebibliography}

\end{document}
